% Part: first-order-logic
% Chapter: introduction
% Section: substitution

\documentclass[../../../include/open-logic-section]{subfiles}

\begin{document}

\olfileid[tr]{fol}{int}{sub}

\olsection{Yerine Koyma}

Her şeyin birbirine nasıl bağlandığını ve sözdizim ile semantiğin
geliştirilmesinin daha sonraki ileri incelemelerimize nasıl temel
hazırladığını göstermek için bir örneği ele alalım. !!{derivation}
sistemlerimiz $\lforall[\Obj v_0][\Atom{\Obj P}{\Obj v_0}]$
ifadesinden $\Atom{\Obj P}{\Obj a}$ ifadesini !!{derive} olanağı
vermelidir. Belki bunu bir çıkarım kuralı olarak bile ifade etmek
isteriz. Fakat bunun için kuralı en genel terimlerle; yalnızca
$\Obj P$, $\Obj a$ ve $\Obj v_0$ için değil, herhangi bir
!!{formula}~$!A$, herhangi bir~$t$ terimi ve !!{variable}~$x$ için
ifade edebilmeliyiz. (!!{constant}s ifadelerinin terim olduğunu
hatırlayın; ayrıca !!{constant}s ve !!{function}s kullanılarak kurulan
daha karmaşık terimleri de ele alacağız.) Dolayısıyla ``ne zaman
$\lforall[x][!A(x)]$ ifadesini !!{derive} ederseniz,~$\lforall[x]$
bölümünü kaldırıp~$x$ yerine~$t$ koyma sonucunda elde edilen~$!A(t)$
ifadesini çıkarmakta haklısınız'' gibi bir şey söyleyebilmek isteriz.
Fakat ``$x$ yerine~$t$ koymak'' tam olarak ne demektir? $!A(x)$ ile
$!A(t)$ arasındaki ilişki nedir? Bu her zaman çalışır mı?

Bunu kesinleştirmek için \emph{yerine koyma} işlemini tanımlarız.
Yerine koyma aslında güçtür; çünkü~$!A$ içindeki bütün~$x$ ifadelerini
öylece~$t$ ile değiştiremeyiz ve her~$t$, herhangi bir~$x$ yerine
konamaz. Bunu yine tümevarımlı tanımlarla ele alacağız. Fakat tanım
yapıldıktan sonra ``$\lforall[x][!A(x)]$ ifadesinden~$!A(t)$ çıkar''
çıkarım kuralı kesin bir tanım olur. Dahası bunun iyi bir çıkarım
kuralı olduğunu, yani $\lforall[x][!A(x)]$ ifadesinin~$!A(t)$
ifadesini mantıksal olarak gerektirdiğini gösterebileceğiz. Ancak bunu
ispatlamak için ilk bakışta ``bunu neden yapıyoruz?'' diye sormanıza
yol açabilecek başka bir şeyi daha ispatlamalıyız.
$\lforall[x][!A(x)]$ ifadesinin~$!A(t)$'yi mantıksal olarak
gerektirmesi, $\Sat M{!A(t)}$ olup olmamasının yalnızca~$t$ teriminin
değerine bağlı olması olgusuna dayanır. Başka bir deyişle $m$,
$\Domain M$ içindeki~$t$ teriminin gösterdiği !!{element} olsun; o
zaman $\Sat M{!A(t)}[s]$ ancak ve ancak
$\Sat M{!A(x)}[\Subst{s}{m}{x}]$ ise. Bu,~$t$ !!{variable}s içerse
bile geçerlidir; fakat sonucu tam olarak nasıl ifade ettiğimize dikkat
etmemiz gerekecektir.

\end{document}
